\chapter{اللقاء الثامن عشر: ختام الجزء الأول من سورة البقرة}

\section{مقدمة}
السلام عليكم ورحمة الله وبركاته. الحمد لله رب العالمين، وأشهد أن لا إله إلا الله وحده لا شريك له، وأشهد أن محمداً عبده ورسوله. اللهم صل على محمد وعلى آل محمد كما صليت على آل إبراهيم إنك حميد مجيد، اللهم بارك على محمد وعلى آل محمد كما باركت على آل إبراهيم إنك حميد مجيد.

نبدأ بإذن الله تبارك وتعالى الدرس السادس عشر (الثامن عشر بترقيم الدورة) من اجتماعنا على قراءة ومدارسة تفسير الإمام \rawi{ابن جرير الطبري} عليه رحمة الله، وارجو أن ننهي اليوم إن شاء الله تبارك وتعالى الجزء الأول (من القرآن الكريم)، وسيكون عندنا لقاء طويل يشارك فيه كل منكم ببحث أو دراسة أو فوائد استخرجها من هذا الكتاب.

\separator

\section{تأويل (وقال الذين لا يعلمون لولا يكلمنا الله أو تأتينا آية)}
\isnad{قال الطبري رحمه الله:} «اختلف أهل التأويل فيمن عنى الله بقوله \textit{وقال الذين لا يعلمون لولا يكلمنا الله}، فقال بعضهم: عنى الله بذلك النصارى (وهو مروي عن \rawi{مجاهد}). وقال آخرون: بل عنى الله بذلك اليهود الذين كانوا في زمان رسول الله ﷺ (وهو مروي عن \rawi{ابن عباس}). وقال آخرون: بل عنى الله بذلك مشركي العرب (وهو مروي عن \rawi{قتادة} و\rawi{الربيع} و\rawi{السدي})».

\isnad{قال أبو جعفر مرجحاً:} «وأولى هذه الأقوال بالصحة والصواب قول القائل إن الله تعالى ذكره عنى بقوله \textit{وقال الذين لا يعلمون} النصارى دون غيرهم. لأن ذلك في سياق خبر الله عنهم وعن افترائهم عليه وادعائهم له ولداً... أما الزاعم أن الله عنى بقوله \textit{وقال الذين لا يعلمون} العرب فإنه قائل قولا لا خبر بصحته ولا برهان على حقيقته في ظاهر الكتاب».

\begin{fawaid}[تقديم ما يشهد له السياق على مجرد الاحتمال]
هذا الموضع من أوضح أمثلة ترجيح \rawi{الطبري} بالسياق؛ فقد كانت الأقوال المحتملة متعددة، لكنه قدم القول الذي ينسجم مع جريان الآيات في خطاب النصارى وذكر افترائهم، ورد غيره لعدم قيام برهان من ظاهر الكتاب عليه. وهذه قاعدة جليلة في التفسير: ليس كل احتمال مذكور سواء، بل أولاها ما شدته قرائن السياق.
\end{fawaid}

\separator

\section{تأويل (كذلك قال الذين من قبلهم مثل قولهم تشابهت قلوبهم)}
\isnad{قال أبو جعفر:} «قد دللنا على أن الذين عنى الله تعالى ذكره بقوله \textit{وقال الذين لا يعلمون لولا يكلمنا الله} هم النصارى، فالذين قالت النصارى مثل قولهم هم اليهود. لأن اليهود سألت موسى عليه السلام أن يريهم ربهم جهرة أو أن يسمعهم كلام ربهم... وكذلك تمنت النصارى على ربها تحكماً... فقلوبهم تشابهت في الكفر بربهم والفرية عليه وتحكمهم على أنبياء الله ورسله...».

\begin{fawaid}[تشابه القلوب يؤدي لتشابه الأفعال وتحديد المشبه به]
نلاحظ أن الطبري رحمه الله يحدد الأقوال بالقرائن. فرجح أن المعني بـ (الذين لا يعلمون) هم النصارى بدلالة السياق، وأن (الذين من قبلهم) هم اليهود، فتشابهت قلوب الطائفتين في العناد. وفي الآية دلالة عظيمة أن تشابه القلوب بالمرض (كالحسد والكبر) ينتج عنه تشابه الأفعال.
\end{fawaid}

\separator

\section{تأويل (قد بينا الآيات لقوم يوقنون)}
\isnad{قال أبو جعفر:} «وخص الله بذلك القوم الذين يوقنون لأنهم هم أهل التثبت في الأمور والطالبون المعرفة على يقين وصحة. فأخبر الله جل ثناؤه أنه بيّن لمن كان هذه الصفة صفته ما بين من ذلك ليزول شكه ويعلم حقيقة الأمر».

\begin{fawaid}[موقف المعاندين من الآيات وواجب الداعية]
الآيات لا ينتفع بها كل أحد، فالمعاند المستكبر لا تنفعه الآيات (ولو جاءتهم كل آية). وهذه الآية تفيد الداعية إلى الله ألا يستجيب لطلبات المعاندين في طلب المزيد من الحجج بعد أن يبين لهم الحق؛ فالاستمرار في الجدال مع المعاند مضيعة للوقت، والأولى صرف الجهد لمن (جاءك يسعى وهو يخشى).
\end{fawaid}

\begin{fawaid}[الفرق بين التسليم واليقين]
الله تبارك وتعالى وصف المؤمنين هنا بـ (يوقنون) ولم يقل (يسلمون)، لأن المؤمن على بينة من ربه، وفرق بين من يسلم أمراً لله وهو غير مستيقن، وبين المسلم المستيقن. فضعيف اليقين يُفتن بأدنى شبهة ويصده المجرمون، فلا بد للمسلم أن يعزز يقينه بقراءة القرآن وسير الأنبياء ليرسخ قلبه وقت الابتلاء.
\end{fawaid}

\separator

\section{تأويل (إنا أرسلناك بالحق بشيرا ونذيرا ولا تسأل عن أصحاب الجحيم)}
قرأ عامة القراء: \textit{وَلَا تُسْأَلُ عَنْ أَصْحَابِ الْجَحِيمِ} بضم التاء ورفع اللام على الخبر. وقرأها بعض أهل المدينة: \textit{وَلَا تَسْأَلْ} جزماً بمعنى النهي.
واستدل أصحاب القراءة الثانية بأحاديث فيها أن النبي ﷺ قال: \begin{hadith}ليت شعري ما فعل أبواي؟\end{hadith} فنزلت الآية.

\isnad{قال أبو جعفر مرجحاً:} «والصواب عندي من القراءة في ذلك قراءة من قرأ بالرفع على الخبر... فبلِّغ رسالتي فليس عليك من أعمال من كفر بك... ولا أنت مسؤول عما عمل بعد ذلك. ولم يجرِ لمسألة رسول الله ﷺ ربه عن أصحاب الجحيم ذكر، فيكون لقوله \textit{ولا تسأل عن أصحاب الجحيم} وجه يوجه إليه... ولا خبر تقوم به الحجة على أن النبي نهي عن أن يسأل في هذه الآية عن أصحاب الجحيم...».

\begin{fawaid}[ترجيح القراءة بالمعنى والسياق ورد الأحاديث الضعيفة]
رجح الطبري قراءة الرفع (الإخبار) على الجزم (النهي) لسببين:
1. ضعف الأحاديث المروية في سبب النزول (بأن النبي سأل عن أبويه). الطبري لا يعتمد على الأخبار المرسلة والضعيفة في الترجيح إذا خالفت السياق.
2. اتساق قراءة الرفع مع سياق الآيات التي تتحدث عن كفر اليهود والنصارى وعنادهم، فكان الأنسب إخبار النبي ﷺ بأنه غير مسؤول عن كفرهم (كقوله: لست عليهم بمصيطر)، وأنه ليس هناك دلالة واضحة على النهي.
\end{fawaid}

وردّ \rawi{الطبري} على \rawi{الأخفش} الذي جعل الآية في موضع "الحال" (أي: غير مسؤول عنهم)، محتجاً بقراءة \rawi{ابن مسعود} (ولن تُسأل) وقراءة \rawi{أبي بن كعب} (وما تُسأل)، فدخول (ولن) و(وما) يمنع أن تكون حالاً، بل هي استئناف.

\separator

\section{تأويل (ولن ترضى عنك اليهود ولا النصارى حتى تتبع ملتهم)}
\isnad{قال أبو جعفر:} «يعني جل ثناؤه... وليست اليهود يا محمد ولا النصارى براضية عنك أبداً، فدع طلب ما يرضيهم ويوافقهم وأقبل على طلب رضا الله في دعائهم إلى ما بعثك الله به من الحق».

\begin{fawaid}[الاطلاع على بواطن الكفار لقطع التعلق بهم]
علّم الله نبيه ﷺ ما في قلوب الكفار والمنافقين ليقطع الطمع في إرضائهم؛ فالنبي ﷺ كان حريصاً على إيمانهم، فأعلمه الله أنهم لن يرضوا إلا باتباع ملتهم، ليتوجه بكليته إلى إرضاء الله. وفيه درس للداعية ألا يداهن أو يسايس على حساب دينه طلباً لرضا المخالفين، بل يبين الحق ولو أغاظهم (قل إن هدى الله هو الهدى).
\end{fawaid}

\separator

\section{تأويل (ولئن اتبعت أهواءهم بعد الذي جاءك من العلم مالك من الله من ولي ولا نصير)}
\isnad{قال أبو جعفر:} «ولئن التمست يا محمد رضا هؤلاء اليهود والنصارى... فصرت من ذلك إلى رضاهم... من بعد الذي جاءك من العلم بضلالتهم... ليس لك يا محمد ولي يلي أمرك... ولا نصير ينصرك من الله».

\begin{fawaid}[أمانة العلم وعقوبة مخالفته]
العلم ميثاق ومسؤولية. فقوله: \textit{بَعْدَ الَّذِي جَاءَكَ مِنَ الْعِلْمِ} يدل على أن العالم يحاسب أشد من غير العالم إذا اتبع هواه. فمن اتبع هواه بعد علمه بالحق فهو متوعد بانقطاع ولاية الله ونصرته. ومخاطبة النبي ﷺ بهذا التهديد هو من باب أولى لأمته.
\end{fawaid}

\separator

\section{تأويل (الذين آتيناهم الكتاب يتلونه حق تلاوته)}
اختلفوا في المعنيين بالآية؛ فقيل هم أصحاب النبي ﷺ، وقيل هم علماء بني إسرائيل الذين آمنوا بمحمد ﷺ وصدقوا التوراة. ورجح \rawi{الطبري} القول الثاني لورود الآية في سياق أخبار أهل الكتاب.

واختلفوا في معنى: \textit{يَتْلُونَهُ حَقَّ تِلَاوَتِهِ}. فذهب \rawi{ابن مسعود} و\rawi{ابن عباس} و\rawi{مجاهد} و\rawi{قتادة} و\rawi{الحسن} إلى أنه: يتبعونه حق اتباعه، يحلون حلاله ويحرمون حرامه ويعملون بمحكمه. وقال آخرون: يقرأونه حق قراءته.
\isnad{ورجح الطبري الأول قائلاً:} «والصواب من القول... أنه بمعنى يتبعونه حق اتباعه... لإجماع الحجة من أهل التأويل على أن ذلك تأويله».

\begin{fawaid}[معنى التلاوة الحقة]
يكاد يغيب المعنى الأصلي للآية (يَتْلُونَهُ حَقَّ تِلَاوَتِهِ) عن كثير من المعاصرين الذين يحصرونها في أحكام التجويد والقراءة. بينما أجمع المفسرون من الصحابة والتابعين على أن "حق التلاوة" هو الفقه والعمل، واتباع ما فيه، وتحليل حلاله وتحريم حرامه. والتلاوة في لغة العرب تأتي بمعنى القراءة وتأتي بمعنى الاتباع (والقمر إذا تلاها).
\end{fawaid}

\begin{fawaid}[منهج الطبري في نقد الأحاديث والاستدلال بها]
من دقة \rawi{الطبري} هنا أنه لا يمرر الأخبار لمجرد اشتهارها في الباب، بل ينظر في أسانيدها ويصرح أحياناً بأن في بعضها نظراً، ثم يبني التفسير على ما يشهد له السياق وإجماع أهل التأويل. وهذا يبين أن كتابه ليس مجرد جمع روايات، بل هو جمع ونقد وترجيح.
\end{fawaid}

\separator

\section{تأويل (أولئك يؤمنون به ومن يكفر به فأولئك هم الخاسرون)}
\isnad{قال أبو جعفر:} «فأخبر الله جل ثناؤه أن المؤمن بالتوراة هو المتبع ما فيها من حلالها وحرامها والعامل بما فيها من فرائض الله... لأن في اتباعها اتباع محمد نبي الله ﷺ وتصديقه».

\begin{fawaid}[الإيمان تصديق واتباع]
صرّح الطبري مراراً أن الإيمان هو التصديق، ولكنه هنا يبين أن هذا التصديق لا يكون نافعاً ولا يسمى إيماناً شرعياً إلا بالاتباع، فقال: "أولئك يؤمنون به" أي الذين عملوا به واتبعوه. فمجرد المعرفة بالكتاب (التي كانت لليهود الكافرين) لا تغني شيئاً بدون العمل.
\end{fawaid}

\separator

\section{تأويل (وإذ ابتلى إبراهيم ربه بكلمات فأتمهن)}
\isnad{قال أبو جعفر:} «وكان اختبار الله تعالى ذكره إبراهيم اختباراً بفرائض فرضها عليه وأمر أمره به، وذلك هو الكلمات».

واختلفوا في صفة هذه الكلمات:
فقيل هي شرائع الإسلام الثلاثون سهماً (وهو مروي عن \rawi{ابن عباس}).
وقيل هي عشر خصال من الفطرة كالختان والمضمضة وقص الشارب (وهو مروي عن \rawi{ابن عباس} و\rawi{قتادة}).
وقيل هي مناسك الحج. وقيل هي ابتلاؤه بالكوكب والقمر والشمس والنار والهجرة والذبح (مروي عن \rawi{الحسن}).

\isnad{ورجح الطبري التوقف عن التعيين لعدم الحجة القاطعة، قائلاً:} «وجائز أن تكون تلك الكلمات جميع ما ذكره من ذكرنا قوله... وجائز أن تكون بعضه... فغير جائز لأحد أن يقول عنى الله بالكلمات... شيئاً من ذلك بعينه دون شيء... إلا بحجة يجب التسليم لها... ولم يصح بشيء من ذلك خبر عن الرسول ﷺ بالنقل».

\begin{fawaid}[منهج الطبري في عدم الترجيح بلا دليل، ونقد المراسيل]
هذا الموضع يبرز منهج الطبري في عدم الجزم بتعيين المبهمات (كالكلمات التي ابتلي بها إبراهيم) طالما لم يرد نص قاطع من القرآن أو السنة الصحيحة، مع احتماله للكل.
كما أورد الطبري حديثين مرفوعين في المسألة وضعّفهما قائلاً: "غير أنهما خبران في أسانيدهما نظر"، مما يؤكد أنه ينقد الأحاديث ولا يسلم بها بمجرد ورودها.
\end{fawaid}

\separator

\section{تأويل (قال إني جاعلك للناس إماما قال ومن ذريتي قال لا ينال عهدي الظالمين)}
سأل إبراهيم عليه السلام ربه أن يجعل من ذريته أئمة مثله، فأجابه الله بأنه لا ينال عهده (الإمامة أو النبوة) الظالمون.
واختلفوا في العهد: فقيل النبوة، وقيل الإمامة، وقيل دين الله، وقيل طاعته، وقيل الأمان في الآخرة.
\isnad{قال الطبري مرجحاً:} «وهو إعلام من الله تعالى ذكره لإبراهيم أن من ولده من يشرك به ويزول عن قصد السبيل ويظلم نفسه وعباده».

\separator

\section{تأويل (وإذ جعلنا البيت مثابة للناس وأمنا)}
\isnad{قال أبو جعفر:} «وأما المثابة... مفعلة من ثاب القوم إلى الموضع إذا رجعوا إليه... فمعنى قوله: مرجعاً للناس ومعاذاً، يأتونه كل عام ويرجعون إليه فلا يقضون منه وطراً».

\isnad{وأما قوله (وأمنا):} «وإنما سماه الله أمناً لأنه كان في الجاهلية معاذاً لمن استعاذ به... فكان الرجل يلقى قاتل أبيه أو أخيه فلا يعرض له».

\separator

\section{تأويل (واتخذوا من مقام إبراهيم مصلى)}
اختلف القراء في (وَاتَّخِذُوا)؛ فقرأها الجمهور بكسر الخاء على الأمر. وقرأها آخرون بفتح الخاء (وَاتَّخَذُوا) على الخبر.
\isnad{رجح الطبري كسر الخاء (الأمر) محتجاً بالحديث:} «للخبر الثابت عن رسول الله ﷺ الذي ذكرناه آنفاً عن عمر بن الخطاب... قلت يا رسول الله لو اتخذت من مقام إبراهيم مصلى، فأنزل الله: واتخذوا...».

ورجح أن (مقام إبراهيم) هو الحجر المعروف الذي في المسجد الحرام، وأن (المصلى) هو مكان الصلاة (وصلاة النبي ﷺ خلفه ركعتين). ورد الأقوال التي وسعت المعنى وجعلته الحج كله أو الحرم كله.

\separator

\section{تأويل (وعهدنا إلى إبراهيم وإسماعيل أن طهرا بيتي للطائفين...)}
\isnad{قال أبو جعفر:} «وامرنا إبراهيم وإسماعيل بتطهير بيته للطائفين... من عبادة الأصنام والأوثان فيه ومن الشرك بالله».

ورجح الطبري أن (للطائفين) هم الدائرون بالبيت، وأن (العاكفين) هم المقيمون والمجاورون بالبيت وإن لم يطوفوا أو يصلوا، وأن (الرُّكَّع السجود) هم المصلون. ففرّق بين معانيها لتأسيس معانٍ جديدة دون التأكيد المكرر.

\separator

\section{تأويل (وإذ قال إبراهيم رب اجعل هذا بلدا آمنا)}
ساق الطبري الأخبار في تحريم مكة؛ فبعضها يذكر أن الله حرّمها يوم خلق السماوات والأرض (حديث ابن عباس)، وبعضها يذكر أن إبراهيم حرّمها (حديث جابر). 
\begin{fawaid}[الجمع بين الأحاديث المتعارضة ظاهرياً]
جمع الطبري بين الحديثين ببراعة قائلاً: «إن الله جل ثناؤه جعل مكة حرماً حين خلقها... بغير تحريم منه لها على لسان أحد من أنبيائه ورسله، ولكن بمنعه جل ثناؤه من أرادها بسوء... فلم يزل ذلك أمرها حتى بوأها الله إبراهيم خليله... فسأل حينئذ إبراهيم ربه إيجاب فرض تحريمها على عباده على لسانه... فأجابه ربه... فصارت محرمة بفرض تحريمها على خلقه على لسان خليله إبراهيم». فالحديثان حق ولا تدافع بينهما.
\end{fawaid}

\separator

\section{تأويل (وارزق أهله من الثمرات من آمن منهم... قال ومن كفر فأمتعه قليلا)}
خصّ إبراهيم عليه السلام المؤمنين بالدعاء بالرزق لأن الله أعلمه قبلها أنه لا ينال عهده الظالمين، فظن إبراهيم أن الرزق الدنيوي كالإمامة لا يُعطى للكفار.
فأجابه الله: \textit{وَمَن كَفَرَ فَأُمَتِّعُهُ قَلِيلًا} (وهي قراءة الجمهور، ورجحها الطبري على قراءة إبراهيم بصيغة الدعاء: فامْتِعْهُ)، أي أن الله يرزق الكافر في الدنيا ليمتعه متاعاً قليلاً إلى وقت مماته، ثم يضطره ويسوقه يوم القيامة إلى عذاب النار.

\separator

\section{تأويل (وإذ يرفع إبراهيم القواعد من البيت وإسماعيل ربنا تقبل منا)}
أورد الطبري الخلاف في قواعد البيت: هل هما أحدثاها أم كانت موجودة من قبل (قواعد بيت آدم أو قبة من ياقوتة)؟
\begin{fawaid}[منهج الطبري في تفاصيل الإسرائيليات الغيبية]
ذكر الطبري الآثار الإسرائيلية في ذلك ثم أعقبها بقاعدة جليلة: «ولا علم عندنا بأي ذلك كان من أي؛ لأن حقيقة ذلك لا تدرك إلا بخبر عن الله أو عن رسوله بالنقل المستفيض... ولا هو مما يدرك علمه بالاستدلال والمقاييس». فهذا يؤسس لمنهج التوقف في الغيبيات التي لم يصح بها نقل عن المعصوم.
\end{fawaid}

\separator

\section{تأويل (ربنا واجعلنا مسلمين لك ومن ذريتنا أمة مسلمة لك وأرنا مناسكنا وتب علينا)}
\isnad{قال أبو جعفر:} «يعنيان بذلك واجعلنا مستسلمين لأمرك خاضعين لطاعتك... وأرنا مناسكنا أي أعلمنا ودلنا عليها».
\isnad{فإن قال قائل:} وهل كانت لهما ذنوب إلى مسألة ربهما التوبة؟
\isnad{قيل:} «إنه لا أحد من خلق الله إلا وله من العمل فيما بينه وبين ربه ما يجب عليه الإنابة منه والتوبة... وجائز أن يكون عنيا بقولهما وتب علينا: وتب على الظلمة من أولادنا وذريتنا».

\begin{fawaid}[عصمة الأنبياء من الذنوب]
قول الطبري هو الصواب، فما من إنسان إلا ويقع في ذنب يحتاج إلى التوبة (كذنب آدم، واستغفار إبراهيم لأبيه)، وهذا لا يقدح في مقام النبوة. فاصطلاح "العصمة من الذنوب مطلقاً" حادث من المتكلمين والرافضة، والأنبياء معصومون في الوحي والتبليغ ولا يُقرون على خطأ بل يوفقون للتوبة الصادقة التي ترفع درجاتهم.
\end{fawaid}

\separator

\section{تأويل (ربنا وابعث فيهم رسولا منهم يتلو عليهم آياتك ويعلمهم الكتاب والحكمة ويزكيهم)}
بيّن الطبري أن هذا الرسول هو محمد ﷺ، فهو دعوة إبراهيم، وأن الحكمة هي: "العلم بأحكام الله التي لا يدرك علمها إلا ببيان الرسول ومعرفة سننه"، وأن التزكية: التطهير من الشرك والإخلاص لله.
\begin{fawaid}[أساس دعوة الإسلام]
لخصت هذه الآية مهمة النبي الخاتم ﷺ، وهي تلاوة الآيات، وتعليم الكتاب (القرآن) والحكمة (السنة والفقه)، وتزكية النفوس وتطهيرها. وكل داعية أسس دعوته على هذا النسق القرآني فقد أفلح، وكل من أهمل تلاوة القرآن وتزكية النفوس به واتجه للمناهج الفلسفية أو الغربية فقد ضل الطريق.
\end{fawaid}

\separator

\section{تأويل (ومن يرغب عن ملة إبراهيم إلا من سفه نفسه)}
\isnad{قال أبو جعفر:} «وأي الناس يذهب عن ملة إبراهيم ويتركها... إلا من سفه نفسه، أي إلا سفيه جاهل بموضع حظ نفسه فيما ينفعها ويضرها في معادها... وملة إبراهيم هي الحنيفية المسلمة».

\begin{fawaid}[أفضلية هدي الأنبياء وفساد مقولات المتفلسفة]
علّق الطبري على اصطفاء الله لإبراهيم مبيناً أن من خالف سنته وسنة محمد ﷺ فهو عدو لله. وهذا يبطل دعوى الفلاسفة وبعض المتكلمين الذين زعموا أن طرقهم العقلية أهدى أو أعمق من طرق الأنبياء؛ فخير الناس هم الرسل، وخير الهدى ما جاؤوا به، وكل من ادعى غير ذلك فقد ضل وخاب.
\end{fawaid}

\separator

\section{تأويل (أم كنتم شهداء إذ حضر يعقوب الموت إذ قال لبنيه ما تعبدون من بعدي)}
هذا توبيخ لليهود والنصارى في دعواهم أن أنبياء الله (كإبراهيم ويعقوب) كانوا على ملتهم. فـ (أم) هنا للاستفهام المنقطع (بمعنى بل أكنتم شهداء؟). 
وقد وصى يعقوب بنيه بالإسلام، \textit{قَالُوا نَعْبُدُ إِلَٰهَكَ وَإِلَٰهَ آبَائِكَ إِبْرَاهِيمَ وَإِسْمَاعِيلَ وَإِسْحَاقَ إِلَٰهًا وَاحِدًا وَنَحْنُ لَهُ مُسْلِمُونَ}. وفي ذكر إسماعيل ضمن "الآباء" مع أنه عم يعقوب دلالة لغوية على أن العرب تطلق لفظ الأب على العم.

\separator

\section{تأويل (تلك أمة قد خلت لها ما كسبت ولكم ما كسبتم)}
\isnad{قال أبو جعفر:} «يقول لليهود والنصارى... دعوا انتحالهم في حال ميلادهم، فإن الدعاوى غير مغنيتكم عند الله شيئاً، إنما يغني عنكم عنده ما سلف لكم من صالح أعمالكم... ولا تُسألون عما كان إبراهيم وإسماعيل... يعملون من الأعمال».

\begin{fawaid}[الاعتبار بالعمل لا بالنسب]
ختم الله تبارك وتعالى هذا الجزء العظيم بوصية قاطعة لأهل الكتاب (ولنا من بعدهم) ألا يتكلوا على أنسابهم أو على فضل أسلافهم من الأنبياء والصالحين؛ فكل امريء يُحاسب بكسبه وعمله. (من بطّأ به عمله لم يُسرع به نسبه)، فالأنساب لا تغني من الحق شيئاً إذا خلا القلب من الإيمان والعمل الصالح.
\end{fawaid}

\separator

بهذا نكون قد أتممنا بحمد الله الجزء الأول من سورة البقرة وتفسير الإمام الطبري. وسنتوقف للمراجعة والمدارسة، نسأل الله أن يبارك في علمنا وعملنا. والسلام عليكم ورحمة الله وبركاته.
